% title: 试题：微局部分析
% renderer: btex
% legacy-title: 用户:Solution/ 试题: 微局部分析
% source: https://bananaspace.org/legacy/wiki/%E7%94%A8%E6%88%B7:Solution/_%E8%AF%95%E9%A2%98:_%E5%BE%AE%E5%B1%80%E9%83%A8%E5%88%86%E6%9E%90
% summary: 从香蕉空间旧站迁移的原始 bTeX 页面。
\setdisplaytitle{试题：微局部分析}
\sethtmltitle{试题：微局部分析}

\section{2025春微局部分析(口试)}
\subsection{同学A}
\begin{enumerate}
    \item[1.]
    \begin{enumerate}
        \item[(1)]算一下球坐标下的Laplace算子.
        \item[(2)]写出Riemann度量下的Laplace算子.
    \end{enumerate}

    \item[2.]局部上拟微分算子是什么形式？大致说一下拟微分算子复合的象征(symbol)怎么算.

    \item[3.]简述椭圆拟微分算子的拟逆的构造过程.
\end{enumerate}
\subsection{同学B}
\begin{enumerate}
    \item[1.]叙述$WF(u)$的定义,其中$u\in\mathcal{D}'(X)$.

    \item[2.]叙述$WF(u)$和$\operatorname{sing}\operatorname{supp}(u)$的关系.

    \item[3.]$WF[I(a,\varphi)]\subset \Lambda_{\varphi}$.

    \item[4.]$WF[K_A(x,y)]=N^{*}_{\Delta}(X\times X)$.

    \item[5.]证明$b(x,\xi)=\mathrm{e}^{-ix\xi}A(\mathrm{e}^{i\Box\xi})$是$A$的symbol.

    \item[6.]``证明'' $b(x,\xi)$关于$a(x,y,\xi)$的渐近展开.
\end{enumerate}